\documentclass[10pt,a4paper]{amsart}
\usepackage[margin=19mm]{geometry}
\usepackage{amsmath,amssymb,mathtools}
\usepackage[scr=rsfs,cal=euler]{mathalfa}
\usepackage{xfrac}
\usepackage[unicode,hidelinks,bookmarks=false]{hyperref}
\newcommand{\ie}{i.e.\ }
\newcommand{\cf}{cf.\ }
\newcommand{\resp}{respectively\ }
\newcommand{\Subsection}{\subsection}
\providecommand{\nobreakdash}{\nobreak\hbox{-}}
\setlength{\emergencystretch}{2em}
\multlinegap0pt
\usepackage{amssymb}
\usepackage{xcolor}
\newtheorem{theorem}{Theorem}
\title{The apolar inner product and a Bombieri type inequality for polynomials}
\author{J. M. Aldaz, A. Bravo, H. Render}
\date{Source: arXiv:2403.10584v1 (CC BY 4.0)}
\begin{document}
\maketitle

\noindent\emph{Source and licence.} The apolar inner product and a Bombieri type inequality for polynomials. Version arXiv:2403.10584v1 (CC BY 4.0), distributed by arXiv under the Creative Commons Attribution 4.0 International licence, http://creativecommons.org/licenses/by/4.0/, as recorded in the arXiv licence metadata for this exact version and shown on the abstract page https://arxiv.org/abs/2403.10584v1. Full licence notice: \texttt{licenses/arxiv-2403.10584-CC-BY-4.0.txt}.\par\smallskip

\noindent\textit{Editorial scope.} Counted unit: the article's theorem BombieriPos (source0.tex lines 312--320). The article's complete original proof of it is delivered with it (source0.tex lines 322--356, one \texttt{proof} environment of 35 lines). No mathematics is written, repaired or re-derived: the statement and the proof are the article's own lines, in the article's own notation, and no retained line is substituted or re-flowed.  No further source-local context is needed: every cross-reference of every retained line resolves inside the two delivered spans.  Nothing was omitted from any retained span, so the package carries no omission record.  No retained line contains a citation, so the wrapper carries no bibliography and no bibliography entry.  No retained line contains a figure environment, an image-inclusion command or drawing code, so no image is delivered and the package declares no asset dependency.  Editorial formatting only, in addition: an AMS \texttt{amsart} preamble that restates the article's own declarations and macros used by the retained lines, namely the environments \texttt{theorem} on separate counters, together with the standard packages those lines need.  The retained environments print in this document as Theorem 1 in order of appearance, whereas the article prints the same items with the numbering of its own full document.  The complete native source of the article as distributed in the arXiv e-print for this exact version is bundled unchanged as \texttt{sources/156184/source0.tex}.
\begin{theorem} \label{BombieriPos}
	Let $P , Q  \in \mathbb{C}[z_{1},\dots,z_{d}]$ be nonzero 
	polynomials with real, non-negative coefficients. Then
	\[
\left\Vert P \cdot Q \right\Vert _{a}
\ge
\left\Vert P \right\Vert _{a}  \left\Vert Q \right\Vert _{a}.
\]		
\end{theorem}

\begin{proof} Let $P$ be a polynomial of total degree $m$, say 
\begin{equation*}
P\left( z \right) =P_{m}\left( z \right) +\cdots +P_{0}\left( z \right),
\end{equation*}
where for each $j=0, \dots ,m,$
 $P_{j}\left( z\right) $ is a homogeneous polynomial
of total degree $j$, and likewise, write
\begin{equation*}
Q\left( z \right) = Q_{n} \left( z \right) +\cdots +Q_{0}\left( z \right).
\end{equation*}
Since the coefficients of $P$ and $Q$ are non-negative, it follows from the definition of the apolar inner product that for all $j,k,r,s$ we always have $\left\langle  P_j Q_k, P_r Q_s\right\rangle_{a} \ge 0$. Thus 
\begin{equation}
\left\langle P Q, P Q\right\rangle_{a} 
=
\left\langle \sum_{j=0}^m
\sum_{k = 0}^{n} P_j Q_k, \sum_{j=0}^m \sum_{k = 0}^{n} P_j Q_k\right\rangle_{a} 
=
\sum_{j=0}^m \sum_{k = 0}^{n} \sum_{r=0}^m
\sum_{s = 0}^{n} \left\langle  P_j Q_k, P_r Q_s\right\rangle_{a} 
\end{equation}
\begin{equation}
\ge
\sum_{j=0}^m \sum_{k = 0}^{n}  \left\langle  P_j Q_k, P_j Q_k\right\rangle_{a} 
=
\sum_{j=0}^m \sum_{k = 0}^{n}  \left\| P_j Q_k\right\|_{a}^2
\ge
\sum_{j=0}^m \sum_{k = 0}^{n}  \left\| P_j \right\|_{a}^2 \left\|Q_k\right\|_{a}^2
\end{equation}
\begin{equation}
=
\sum_{j=0}^m   \left\| P_j \right\|_{a}^2 \sum_{k = 0}^{n} \left\|Q_k\right\|_{a}^2
= 
\left\| P \right\|_{a}^2 \left\| Q \right\|_{a}^2.
\end{equation}
\end{proof}

\end{document}
